{\huge\chapter{Теоретические основы реидентификации объектов}}
\label{sec:Theory} \index{Theory}

\section{Определения и постановка задачи}
\textbf{Реидентификация (Re-identification, ReID)} — задача поиска изображений одного и того же физического объекта (в работе — транспортного средства) в \emph{галерее} \(\mathcal{G}\) по \emph{запросу} \(q\), когда наблюдения получены с разнесённых по пространству и времени камер, под различными ракурсами и в различных условиях освещения. Модель должна строить отображение \(f_\theta: \mathcal{X}\!\to\!\mathbb{R}^d\), переводя изображение объекта \(x\) в компактный дескриптор \(z=f_\theta(x)\) так, чтобы в метрике \(D(\cdot,\cdot)\) расстояния между экземплярами одного ID были меньше, чем между различными ID:
\begin{equation}
D\big(f_\theta(x_i), f_\theta(x_j)\big) \ll D\big(f_\theta(x_i), f_\theta(x_k)\big),\;\; \text{если } id(x_i)=id(x_j)\neq id(x_k).
\end{equation}
Ключевая особенность ReID — \emph{невидимость идентичностей на обучении}: множества ID в обучающем и тестовом подвыборках не пересекаются, т.е. требуется обобщение на новые экземпляры.

\section{Инвариантность и вариативность наблюдений}
Теоретически качество ReID определяется балансом между \emph{инвариантностью} признаков к факторам съёмки и \emph{дискриминативностью} по отношению к различным экземплярам. Для транспортных средств важны следующие факторы вариативности:
\begin{itemize}
  \item \textbf{Геометрия и ракурс:} вид спереди/сзади/сбоку, перспектива, масштаб;
  \item \textbf{Фотометрия:} освещение (день/ночь), тени, блики, баланс белого;
  \item \textbf{Сенсорика:} различия камер (матрица, оптика, компрессия, шум);
  \item \textbf{Сцена:} частичная окклюзия, динамический фон, погодные условия.
\end{itemize}
Модель должна «подавлять» вариативность, не связанную с идентичностью (style), и усиливать признаки, устойчиво отличающие конкретный экземпляр (shape/parts/texture).

\section{Эмбеддинги и метрическое пространство}
Пусть \(z=f_\theta(x)\in\mathbb{R}^d\) — эмбеддинг изображения. На практике используются L2-нормированные представления \(\hat z = z/\|z\|_2\), что делает \emph{косинусное расстояние} эквивалентным евклидову. Сопоставление выполняется по \(D(\hat z_i,\hat z_j)\in[0,2]\) (евклидово) или по \(\mathrm{sim}(\hat z_i,\hat z_j)=\hat z_i^\top \hat z_j\in[-1,1]\) (косинусная близость). Нормализация улучшает устойчивость к масштабу признаков и стабилизирует обучение при комбинировании функций потерь.

\section{Функции потерь для обучения дискриминативных признаков}
\subsection{Кросс-энтропия с ослаблением меток}
Классификационная ветвь рассматривает каждый идентификатор как отдельный класс. Для выхода классификатора \(p(y|x)\) минимизируется
\begin{equation}
\mathcal{L}_{CE} = - \sum_{c} y^{(\varepsilon)}_c \log p(c|x), \quad
y^{(\varepsilon)}_c = (1-\varepsilon)\,\mathbf{1}\{c{=}y\} + \frac{\varepsilon}{C},
\end{equation}
где \(\varepsilon\in[0,1]\) — параметр label smoothing, \(C\) — число классов (ID в батче/эпохе). Такая регуляризация уменьшает переобучение на «жёсткие» метки и повышает обобщающую способность.

\subsection{Контрастивное и триплетное обучение}
\paragraph{Контрастивная потеря.} Для пары \((x_i,x_j)\) с меткой \(y_{ij}\in\{0,1\}\) минимизируется
\(\mathcal{L}_{contr} = y_{ij}\,D^2 + (1-y_{ij})\max(0,m - D)^2\), где \(m>0\) — зазор.
\paragraph{Triplet Loss.} Для триплета \((a,p,n)\) (якорь, позитив, негатив):
\begin{equation}
\mathcal{L}_{triplet} = \max\left(0, m + D(\hat z_a,\hat z_p) - D(\hat z_a,\hat z_n)\right).
\end{equation}
Эффективность зависит от стратегии семплинга: в ReID распространён батчевый \emph{PK-семплинг} (P идентичностей, K изображений каждой) и \emph{hard-example mining} внутри батча (выбор самых «трудных» позитивов/негативов).

\subsection{Круговые/центрические потери}
Альтернативами являются Circle loss, ArcFace/AM-Softmax и Center loss, которые вводят угловые/аддитивные зазоры в пространстве признаков, улучшая разделимость кластеров идентичностей. Их можно комбинировать с кросс-энтропией или триплетом.

\section{Архитектурные элементы и нормализация признаков}
\paragraph{BNNeck.} Практически значимый приём — отделение «до-BN» представления для триплет-потери и «после-BN» для кросс-энтропии, что уменьшает конфликт целей и стабилизирует оптимизацию. В инференсе обычно используют L2-нормированное «после-BN» представление. % \cite{Luo2019StrongBaseline}
\paragraph{Part-based и attention-модули.} Разделение на регионы (PCB/Aligned-ReID) и механизмы внимания фокусируют модель на дискриминативных деталях автомобиля (форма фар, крыша, диски), повышая устойчивость к ракурсу и окклюзиям.
\paragraph{Transformer-бэкбоны.} Глобальное внимание (ViT-подобные архитектуры) улучшает согласованность признаков на сложных сценах; требует аккуратных аугментаций и регуляризации.

\section{Постобработка ранжирования}
\paragraph{k-reciprocal re-ranking.} Улучшает порядок результатов, учитывая взаимность соседства: если \(x_i\) входит в \(k\)-NN множества \(x_j\) \emph{и наоборот}, то пара считается более надёжной. Итоговая мера схожести строится из сочетания евклидовой/косинусной близости и джаккардовой по расширенным окрестностям. Метод заметно повышает mAP при небольших вычислительных расходах. % \cite{Zhong2017Rerank}

\section{Масштабирование поиска и индексация}
При больших галереях линейный перебор неэффективен. Используются структуры \emph{приближённого ближайшего соседа} (ANN), например, \texttt{FAISS} с IVF/PQ/HNSW, которые ускоряют поиск при умеренной потере точности. Дополнительно применяются сжатие эмбеддингов (PQ/OPQ) и off-line индексация.

\section{Доменный сдвиг и камера-специфические эффекты}
Разные камеры обладают различной цветометрией и шумовыми характеристиками; сцены отличаются погодой и временем суток. Для повышения переносимости применяются:
\begin{itemize}
  \item \textbf{Аугментации стиля} (Color Jitter, Random Erasing, MixStyle/CamStyle);
  \item \textbf{Нормализации}, снижающие зависимость от статистик камеры (InstanceNorm/IBN);
  \item \textbf{Unsupervised/Domain Adaptation} — дообучение на целевом домене без меток (псевдометки, кластеризация, self-training).
\end{itemize}

\section{Оценивание качества}
Качество оценивается метриками \emph{CMC@k} и \emph{mAP}. Для запроса \(q\) с множеством релевантов \(\mathcal{R}_q\) \emph{Average Precision} определяется как
\begin{equation}
AP(q) = \frac{1}{|\mathcal{R}_q|} \sum_{n=1}^{N} P@n \cdot \mathbf{1}\{r_n \in \mathcal{R}_q\}, \quad
mAP = \frac{1}{|\mathcal{Q}|} \sum_{q \in \mathcal{Q}} AP(q).
\end{equation}
CMC@k сообщает долю запросов, для которых хотя бы одно правильное соответствие попало в топ-\(k\) ранжирования.

\section{Этические и правовые аспекты}
В работе исключается использование распознавания регистрационных знаков; идентификация выполняется по внешнему виду и предназначена для агрегированного анализа трафика и инженерных приложений. Следует учитывать локальные нормативы о защите персональных данных и приватности видеонаблюдения; при публикации результатов необходимо анонимизировать примеры (кропы), исключая потенциально идентифицирующие атрибуты.

\section{Выводы по главе}
Сформулированы формальная постановка ReID и теоретические основы построения устойчивых эмбеддингов: нормализация, комбинированные функции потерь, архитектурные приёмы (BNNeck, attention/parts), постобработка (re-ranking) и методы повышения переносимости к доменному сдвигу. 

Эти положения определяют выбор моделей и протоколов, применяемых далее 

в работе.


